%% ==========================================================================
%% related-work.tex — Related Work
%% Paper: "From Taxonomy to Testbed: Quantifying Environmental Attacks
%%         on Multi-Agent A2A Systems"
%% ==========================================================================

\section{Related Work}
\label{sec:related}

Our work bridges several distinct research threads: prompt injection
and content manipulation attacks, classical adversarial machine learning,
multi-agent system security, RAG and knowledge-base poisoning, and
LLM safety evaluation benchmarks.  Table~\ref{tab:related-positioning}
summarizes how our testbed complements and extends prior work.

% ------------------------------------------------------------------
\subsection{Prompt Injection \& Content Manipulation}
\label{sec:rw-injection}

The foundational prompt injection threat was formalized by
Perez and Ribeiro~\cite{perez2022prompt}, who demonstrated that
adversarial prompts can override system-level instructions.
Greshake et~al.~\cite{greshake2023indirect} introduced the more
insidious \emph{indirect} variant: adversarial payloads planted in
web pages, emails, and documents that agents encounter during
tool-augmented retrieval.  Their work showed that LLM-integrated
applications could be compromised without any direct user interaction,
motivating our Content Injection trap category (\S\ref{sec:results-ci}).

Liu et~al.~\cite{liu2024formalizing} formalized prompt injection
against LLM-integrated applications, providing an attack taxonomy and
demonstrating injection through multiple channels.
Yi et~al.~\cite{yi2023benchmarking} contributed a benchmark for
indirect prompt injection defense evaluation.
Zhan et~al.~\cite{zhan2024injecagent} developed \textsc{InjecAgent},
benchmarking indirect prompt injections specifically in
tool-integrated LLM agents.

Our work differs from these contributions in two key ways.
First, we go beyond single-turn injection to study \emph{environmental}
attacks where the adversarial payload is embedded in the agent's
operating environment (CSS-invisible text, HTML comments, image
metadata, dynamically cloaked content) rather than in the user
prompt.  Second, we measure injection success across a controlled
matrix of 2~conditions $\times$ 5~repetitions
with statistical rigor (Bonferroni-corrected $p$-values, Cohen's $d$
effect sizes), providing the quantitative grounding that prior work
has lacked.

% ------------------------------------------------------------------
\subsection{Adversarial Machine Learning Foundations}
\label{sec:rw-adversarial}

The adversarial examples paradigm was established by
Goodfellow et~al.~\cite{goodfellow2015adversarial}, who showed that
small input perturbations could reliably fool neural classifiers.
Carlini and Wagner~\cite{carlini2017evaluating} raised the
evaluation bar by proposing optimization-based attacks that
defeated most known defenses, establishing that defensive
claims must be backed by strong adaptive evaluations.
Madry et~al.~\cite{madry2018towards} introduced adversarial
training as a principled defense, while
Biggio et~al.~\cite{biggio2013evasion} formalized evasion attacks
at test time for traditional ML pipelines.

In the LLM era, Zou et~al.~\cite{zou2023universal} demonstrated
universal adversarial suffixes that transfer across aligned language
models, while Wei et~al.~\cite{wei2024jailbroken} systematically
analyzed how LLM safety training fails through competing
objectives and mismatched generalization.
Chao et~al.~\cite{chao2023jailbreaking} showed that black-box
jailbreaks can be automated in as few as twenty queries.
Liu et~al.~\cite{liu2024autodan} developed \textsc{AutoDAN}, which
generates stealthy jailbreak prompts automatically.

Our testbed extends this adversarial perspective from model-level
perturbations to \emph{environment-level} attacks.  While prior work
focuses on crafting inputs that bypass alignment at the model
boundary, we study how adversarial content in the agent's operating
environment---web pages, documents, API responses, and inter-agent
messages---can manipulate behavior without ever touching the system
prompt or the user's direct input.  This distinction is critical
for multi-agent A2A systems where agents consume content from
diverse, untrusted sources.

% ------------------------------------------------------------------
\subsection{Multi-Agent Security}
\label{sec:rw-multiagent}

The emerging multi-agent LLM ecosystem introduces attack surfaces
that go beyond single-agent threats.  Cohen et~al.~\cite{cohen2024unleashing}
demonstrated self-replicating adversarial attacks that propagate
across multi-agent GenAI systems, compromising new agents through
inter-agent communication---a ``worm'' for LLM-based systems.
Gu et~al.~\cite{gu2024agent} showed that a single adversarial
image can exponentially compromise a population of multimodal
agents through viral propagation, with attack success scaling
with population size.
Tian et~al.~\cite{tian2024evil} explored safety risks specific
to LLM-based agent systems, revealing emergent dangerous behaviors
when multiple agents collaborate.

The A2A protocol~\cite{google2025a2a} defines standardized
inter-agent communication, enabling multi-agent orchestration through
agent cards~\cite{google2025agentcard}, task delegation, and
artifact exchange.  Frameworks such as
AutoGen~\cite{wu2023autogen} and MetaGPT~\cite{hong2024metagpt}
demonstrate the practical viability of multi-agent collaboration
but do not systematically evaluate adversarial robustness.

Our work is the first to evaluate environmental attacks
\emph{specifically within the A2A protocol context}.  We test
cascade failure propagation (Systemic traps, \S\ref{sec:results-sys}),
where a compromised agent poisons messages flowing through
the A2A task graph.  We also introduce agent impersonation and
message poisoning scenarios unique to protocol-mediated
multi-agent systems.  Chen et~al.~\cite{chen2024agentpoison}
concurrently studied agent poisoning through memory and knowledge
bases, which complements our environment-level attack model.

% ------------------------------------------------------------------
\subsection{RAG \& Knowledge Poisoning}
\label{sec:rw-rag}

Retrieval-Augmented Generation (RAG)~\cite{lewis2020rag} enables
LLMs to access external knowledge, but also opens a new attack
surface: the retrieval corpus itself.
Zhong et~al.~\cite{zhong2023poisoning} demonstrated corpus
poisoning by injecting adversarial passages that manipulate
retrieval rankings and downstream generation.
Zou et~al.~\cite{zou2024poisonedrag} developed \textsc{PoisonedRAG},
showing that even a small number of poisoned documents can
corrupt RAG output with high success rates.
Xue et~al.~\cite{xue2024badrag} identified vulnerabilities
in RAG pipelines through targeted poisoning attacks that
evade existing detection mechanisms.
Chaudhari et~al.~\cite{chaudhari2024phantom} proposed
\textsc{Phantom}, demonstrating general-purpose trigger attacks
on retrieval-augmented generation that activate only in the
presence of specific query patterns.

Our Cognitive State trap category (\S\ref{sec:results-cs})
extends RAG poisoning to the multi-agent context.  We implement
four scenarios---vector embedding poisoning, gradual knowledge
drift, cross-contamination between agents, and ranking
manipulation---that target a live RAG pipeline via the A2A
protocol.  Unlike prior work that evaluates poisoning in
isolated single-model settings, we measure how poisoned knowledge
propagates through agent-to-agent communication, amplifying
factual degradation across the system.

% ------------------------------------------------------------------
\subsection{LLM Safety Benchmarks}
\label{sec:rw-benchmarks}

Several benchmarks evaluate LLM safety and trustworthiness.
Mazeika et~al.~\cite{mazeika2024harmbench} proposed
\textsc{HarmBench}, a standardized framework for automated red
teaming with robust refusal evaluation.
Wang et~al.~\cite{wang2024decodingtrust} developed
\textsc{DecodingTrust}, a comprehensive GPT trustworthiness
assessment covering toxicity, bias, robustness, and privacy.
Zhang et~al.~\cite{zhang2024safetybench} introduced
\textsc{SafetyBench}, evaluating LLM safety via multiple-choice
questions across diverse risk categories.
Sun et~al.~\cite{sun2024trustllm} proposed \textsc{TrustLLM}
for holistic trustworthiness evaluation.
Debenedetti et~al.~\cite{debenedetti2024agentdojo} developed
\textsc{AgentDojo}, specifically evaluating adversarial robustness
of LLM agents in tool-use scenarios.
Ruan et~al.~\cite{ruan2024identifying} proposed an LM-emulated
sandbox for identifying risks in LM agents.

These benchmarks focus predominantly on single-model or
single-agent evaluation.  Our testbed is distinguished by three
design choices: (1)~it evaluates \emph{environmental} rather than
\emph{prompt-level} attacks, implementing the full taxonomy from
Yuntao et~al.~\cite{yuntao2025agenttraps}; (2)~it tests attacks
within a real multi-agent A2A protocol system rather than in
isolated model calls; and (3)~it provides a controlled experimental
matrix (22~scenarios $\times$ 2~conditions
$\times$ 5~reps $=$ 220 main runs) with non-parametric statistical tests
(Wilcoxon signed-rank~\cite{wilcoxon1945individual}) and
calibrated effect sizes (Cohen's $d$~\cite{cohen1988statistical}).

% ------------------------------------------------------------------
\subsection{Positioning Summary}

Table~\ref{tab:related-positioning} compares our testbed against
the most closely related prior work across five evaluation dimensions.

\begin{table}[t]
\centering
\caption{Positioning against prior work.
  \ding{51} = addressed; \ding{55} = not addressed; $\circ$ = partially.}
\label{tab:related-positioning}
\small
\begin{tabular}{@{}lccccc@{}}
\toprule
\textbf{Work} &
  \rotatebox{70}{\textbf{Environmental}} &
  \rotatebox{70}{\textbf{Multi-Agent}} &
  \rotatebox{70}{\textbf{A2A Protocol}} &
  \rotatebox{70}{\textbf{Full Taxonomy}} &
  \rotatebox{70}{\textbf{Stat.\ Rigor}} \\
\midrule
Greshake et al.~\cite{greshake2023indirect}
  & $\circ$ & \ding{55} & \ding{55} & \ding{55} & \ding{55} \\
Perez \& Ribeiro~\cite{perez2022prompt}
  & \ding{55} & \ding{55} & \ding{55} & \ding{55} & \ding{55} \\
Zou et al.~\cite{zou2023universal}
  & \ding{55} & \ding{55} & \ding{55} & \ding{55} & $\circ$ \\
Cohen et al.~\cite{cohen2024unleashing}
  & $\circ$ & \ding{51} & \ding{55} & \ding{55} & \ding{55} \\
Mazeika et al.~\cite{mazeika2024harmbench}
  & \ding{55} & \ding{55} & \ding{55} & \ding{55} & \ding{51} \\
Debenedetti et al.~\cite{debenedetti2024agentdojo}
  & $\circ$ & \ding{55} & \ding{55} & \ding{55} & $\circ$ \\
Yuntao et al.~\cite{yuntao2025agenttraps}
  & \ding{51} & $\circ$ & \ding{55} & \ding{51} & \ding{55} \\
\midrule
\textbf{This work}
  & \ding{51} & \ding{51} & \ding{51} & \ding{51} & \ding{51} \\
\bottomrule
\end{tabular}
\end{table}

The Yuntao et~al.\ taxonomy~\cite{yuntao2025agenttraps} provides the
conceptual foundation; our contribution is to operationalize it as a
reproducible, statistically rigorous testbed and to extend it to the
multi-agent A2A protocol setting.  We further contribute seven
purpose-built mitigations and empirical evidence of both their
effectiveness and their failure modes (including iatrogenic regressions).
